\documentclass[12pt]{article}
%All the packages the project needs
%--------------------------------------------
%Margins
\usepackage[top=1.5cm,
bottom=1.5cm,
left=2.5cm,
right=2cm,
headsep=15pt,
headheight=25pt,
letterpaper,
includehead,includefoot]{geometry} 
%-------------------------------------------

%-------------------------------------------
%Codification and language
%Document language
\usepackage[spanish,es-noshorthands]{babel}
% Required for including letters with accents
\usepackage[utf8]{inputenc} 
% Use 8-bit encoding that has 256 glyphs
\usepackage[T1]{fontenc}
%-------------------------------------------

%-------------------------------------------
%Colours
\usepackage{color}
\usepackage{xcolor}
%-------------------------------------------

%-------------------------------------------
%Fonts and text
%Add a new Font
\usepackage{addfont}
\addfont{OT1}{rsfs10}{\rsfs}
%Allows to modify the font of a section
\usepackage{sectsty}
%Latin Modern fonts - Allows to resize the font
\usepackage{lmodern}
%Allows to write code
\usepackage{verbatim}
%Dumb text
\usepackage{lipsum}
%Include currency symbols
\usepackage{textcomp} 
%-------------------------------------------

%-------------------------------------------
%Pictures and graphics
%Allows to include images
\usepackage{graphicx}
\usepackage{sidecap}
%Most powerful image package
\usepackage{tikz}
\usetikzlibrary{babel}
\usetikzlibrary{positioning}
%Sometimes used images packages

\usepackage{pst-all}
\usepackage{pstricks}
%Allows to position figures inside a paragraph
\usepackage{wrapfig}
%Information needed
%https://en.wikibooks.org/wiki/LaTeX/Floats,_Figures_and_Captions
%-------------------------------------------
%-------------------------------------------
%Math packages
\usepackage{mathtools}
%\usepackage{unicode-math}
\usepackage{amsmath}
\usepackage{amssymb}
\usepackage{amsfonts}
%Nice cancelations arrows
\usepackage{cancel}
%Vector arrows on a character
\usepackage{esvect}
%-------------------------------------------

%-------------------------------------------
%Table and listing packages
\usepackage{multirow}
\usepackage{array}
\usepackage{diagbox}
\usepackage{listings}
\usepackage[shortlabels]{enumitem} 
%Differents images and tables in the same space
\usepackage{subcaption}
%Use to places the float at precisely the location in the LaTeX code´and other stuff
\usepackage{float}
%Flexible tables
\usepackage{tabu}
\usepackage{multicol}
%-------------------------------------------

%-------------------------------------------
%Refs and cites
%\usepackage{achemso}
%\usepackage{natbib}
\usepackage{hyperref}
\hypersetup{
colorlinks=true,
urlcolor= blue,
citecolor=blue,
linkcolor= black,
bookmarks=true,
bookmarksopen=true,
}
\usepackage{csquotes}
\usepackage[backend=biber]{biblatex}
%\usepackage{cite}
%-------------------------------------------

%-------------------------------------------
%Boxes
\usepackage{tcolorbox}
\usepackage{fancybox}
\usepackage{colortbl}
%-------------------------------------------

%-------------------------------------------
%Page style
\usepackage[usestackEOL]{stackengine}
\usepackage{fancyhdr}
\pagestyle{fancy}
\newsavebox{\UNl}
\newsavebox{\UNd}
\setlength{\headheight}{2.5cm}
\lhead{\usebox{\UNl}\usebox{\UNd}}
\rhead{\slshape Proyecto de Tesis}
\cfoot{}
\rfoot{\thepage}
%\lfoot{\today}
\renewcommand{\headrulewidth}{1pt}

%-------------------------------------------

%-------------------------------------------
%others
\usepackage{pgfplots,caption}
\usepackage{comment}
\usepackage{pdflscape}
\usepackage{rotating}
\usepackage{parskip}
\usepackage{hhline}
\usepackage{longtable,lscape}
\usepackage{booktabs,bigstrut}
\usepackage{enumerate}
\usepackage{cleveref}
\usepackage{braket}

\addbibresource{bibliografia.bib}
\title{Proyecto de Tesis}
\author{Eduardo Andres Delgadillo Monsalve}
\date{\today}
%Inicio del documento
\pgfplotsset{compat=1.17}

\begin{document}
\renewcommand{\tablename}{Tabla}%% Para cambiar "Cuadro" a "Tabla"
\renewcommand{\listtablename}{Índice de tablas}
%Inicio de la portada
\begin{titlepage}
    \centering
    \thispagestyle{empty}
    \begin{center}
        \begin{figure}
        \centering%
        \includegraphics{images/EscudoUN.png}
    \end{figure}
    
    \vspace{3cm}
    
    \textbf{CARACTERIZACIÓN DE LAS PROPIEDADES DE LA REGOLITA LUNAR MEDIANTE EL DISEÑO E IMPLEMENTACIÓN DE UN MODELO ANÁLOGO MORFOLÓGICO Y UN MODELO ESPECTRAL}\\[1in]    
    Eduardo Andrés Delgadillo Monsalve \\  [3in]

   \textbf{Universidad Nacional de Colombia}\\
   Facultad de Ciencias\\
   Bogotá D.C.\\
   \today
    \end{center}
\end{titlepage}

\newpage
\begin{titlepage}
    \thispagestyle{empty}
    \begin{center}
        \begin{figure}
        \centering%
        \includegraphics{images/EscudoUN.png}
    \end{figure}
    
    \vspace{1cm}
    
        \textbf{CARACTERIZACIÓN DE LAS PROPIEDADES DE LA REGOLITA LUNAR MEDIANTE EL DISEÑO E IMPLEMENTACIÓN DE UN MODELO ANÁLOGO MORFOLÓGICO Y UN MODELO ESPECTRAL}\\[1in]    
    Eduardo Andrés Delgadillo Monsalve \\  [1.5 cm]
 Proyecto de tesis presentado como requisito parcial paara optar al titulo de:\\[3mm] \textbf {\large{Magister en Ciencias Astronomia}}\\[1cm]
 Director:\\[3mm] Mario Armando Higuera Garzon\\[1mm] Observatorio Astronómico Nacional\\[1cm]
 Codirector:\\[3mm] David Ardila\\ Jet Propulsion Laboratory\\[1in]
 \textbf{Universidad Nacional de Colombia  }\\
   Facultad de Ciencias\\
   Bogotá D.C.\\
   \today
    \end{center}
\end{titlepage}

\newpage
\thispagestyle{empty}
\tableofcontents
\thispagestyle{empty}
\newpage

% \listoffigures

% \listoftables

\newpage

\setcounter{page}{1}
\section{Introducción}
La Luna ha sido uno de los objetos mas estudiados y de mayor interés para la humanidad desde épocas antiguas. En la mayoría de civilizaciones y culturas,
de las que se tienen registro, la Luna ha tenido un papel fundamental dentro de su cosmogonía. Este satélite natural se considero por mucho tiempo como 
la contraparte de la estrella que domina el cielo diurno, el Sol. Eventos naturales como los eclipses de Sol, los eclipses de Luna y el ciclo lunar, fueron interpretados
y estudiados a lo largo de toda la historia, incluso uno de los mayores logros como especie humana ha sido el alunizaje de humanos sobre la superficie de la Luna, 
en el año 1969. Desde entonces, el estudio de la Luna se ha incrementado y representa una de las ramas de investigación mas importantes en la búsqueda por habitar su superficie.

Debido al proceso de formación y evolución, la superficie de la Luna esta compuesta principalmente de polvo y regolito. Las diferentes misiones de exploración lunares de las distintas 
agencias espaciales han recopilado muestras que han sido analizadas y caracterizadas. Aunque algunas propiedades de este material se conocen bien, algunas otras y que son muy importantes 
requieren aún de una investigación mas profunda. Debido a la dificultad de obtener dichas muestras de este material, técnicas como la reproducción o la implementación de un modelo análogo 
que replique sus propiedades son necesarias. La caracterización del material de la superficie lunar puede hacerse utilizando su morfología, es decir su composición, 
forma, tamaño, estructura y estado, también utilizando su emisión térmica, que presenta variaciones debido a las diferencias de temperatura que se pueden presentar por su ciclo 
y a su abrupto relieve, y por ultimo utilizando la distribución espectral de energía (SED) del material de la superficie, principalmente los granos de polvo.

El conocimiento completo de las propiedades del material de la superficie de la Luna es de gran importancia sobre todo para posteriores misiones de alunizaje. 
El conjunto de todas estas propiedades definirán las características del lanzamiento y alunizaje de cualquier misión de exploración.
\section{Objetivos}
\subsection{Objetivo General}
Caracterizar las propiedades de la regolita lunar mediante el diseño e  implementación de un modelo análogo morfológico y un modelo espectral.

\subsection{Objetivos específicos}
\begin{itemize}
    \item Estudiar y caracterizar las propiedades de la regolita  en la superficie lunar
    \item Diseñar un modelo análogo morfológico que permita reproducir las condiciones en la superficie lunar %Para la region del crater Garavito Armero
    \item Diseñar un modelo análogo de curva de emisión que permita comparar el SED del grano de la superficie lunar con el SED obtenido con la herramienta CIGALE
\end{itemize}

\section{Antecedentes}
Artículos, estudios, análisis preliminares existentes sobre este tema
\subsection{Estado del Arte}
Todo el background o marco teórico Según Einstein
\section{Planteamiento del problema}
Definir la motivación para el proyecto Según Einstein \cite{watson53}.

\printbibliography
\end{document}